\begin{definition}\label{4.1.1}
	Let $X$ be a topological space. The \emph{singular set} of $X$ is the simplicial set $\Sing_{\bullet} (X)$ given by
	\begin{itemize}
		\item $\Sing_n(X)=\Top (\left| \Delta^n \right| ,X)$;
		\item Degeneracies and faces are given by precomposition:
			\[\begin{tikzcd}[row sep = 0em, column sep = 2.5em]
			d_i : \Top (\left| \Delta^n \right| ,X)\ar[" (d^i)^* ", r]  & \Top (\left| \Delta^{n-1}  \right| ,X) \\
			s_i : \Top (\left| \Delta^n \right| ,X)\ar[" (s^i)^* "', r]  & \Top (\left| \Delta^{n+1}  \right| ,X).
			\end{tikzcd}\]
	\end{itemize}
\end{definition}

\begin{remark}\label{4.1.2}
	It is immediately checked that this is a simplicial set since $\Hom(-, x) $ is contravariant for all $x$ in all categories $\mathcal{C} $, and since the geometric cofaces and codegeneracies satsify the cosimplicial identities, precomposition must satisfy the faces and degeneracies. One may also view this as a consequence of Yoneda's lemma.
\end{remark}

\begin{proposition}\label{4.1.3}
	The singular complex is functorial $\Sing : \Top  \to \sSet $. Action on morphisms is given by levelwise postcomposition.
\end{proposition}
\begin{proof}
	The singular complex is functorial levelwise, meaning for each $n$, there is a functor $\Sing_n : \Top  \to \Set $ given by $\Sing_n = \Top (\left| \Delta^n \right| ,-)$. One might hope that these assemble into a functor $\Top \to \sSet $. Note that for any $f : X \to Y$ in $\Top $, the map induced at each level is simply postcomposition $\Top(\left| \Delta^n \right| ,f)=f_*$. It's well-known that $(fg)_*=f_*g_*$ and that $1_*$ is the identity by functoriality of $\Hom(x,-)$, so we need only show $f_*$ is a morphism of simplicial sets. By \cref{3.1.5}, we must show the diagrams
	\[\begin{tikzcd}[row sep = 2.5em, column sep = 2.5em]
		\Top (\left| \Delta^x \right| ,X)\ar[" f_* ", r] \ar[" (d^i)^* "', d]  & \Top (\left| \Delta^n \right| ,Y)\ar[" (s^i)^* ", d]  & \Top (\left| \Delta^n \right| ,X)\ar[" f_* ", r] \ar[" (s^i)^* "', d]  & \Top (\left| \Delta^{n}  \right| ,Y)\ar[" (d^i)^* ", d]  \\
	\Top (\left| \Delta^{n-1}  \right| ,X)\ar[" f_* "', r]  & \Top (\left| \Delta^{n-1}  \right| ,Y) & \Top (\left| \Delta^{n+1}  \right| ,X)\ar[" f_* "', r]  & \Top (\left| \Delta^{n+1}  \right| ,Y)
	\end{tikzcd}\]
	commute. These diagrams boil down to the equalities
	\[
		f_*(d^i)^*=(d^i)^*f_*,\qquad f_*(s^i)^*=(s^i)^*f_*,
	\]
	which are easily verified by tracing a map $g$ around the squares, and noting that we are precomposing by the same degeneracies and faces in both verticals of each square.
\end{proof}

\begin{notation}\label{4.1.4}
	Many authors prefer to write $S(X)$ for the singular set, and $S_n(X)$ for its set of $n$-simplices. We prefer writing $\Sing$, since we think it is less ambiguous and because we think it looks cooler.
\end{notation}


The singular complex functor is of great importance for a variety of reasons. For example, singular sets are always \emph{Kan complexes}, which have significant homotopic importance and are a model for $\infty $-groupoids; there is an adjunction $|-|\dashv \Sing$ which is a \emph{Quillen equivalence}; and in particular the geometric realization of a singular set is homotopy equivalent to the original topological space. For now, we explore a more simple application of the singular complex functor which also has extreme importance in many fields: the singular homology of a topological space.

\begin{remark}\label{4.1.5}
	In \cref{4.1.6}, we will consider \emph{simplicial abelian groups}, which is a contravariant functor $X : \boldsymbol{\Delta} ^{\mathrm{op}}  \to \Ab$. More explicitly, they are groups $X_n$ for each $n\geq 0$, and \emph{group homomorphism} $d_i : X_n \to X_{n-1} $ and $s_i : X_n \to X_{n+1} $ satisfying the simplicial identities (\cref{3.1.7}). This definition is precisely the same as a simplicial set, but using $\Ab$ instead of $\Set $. One can equivalently define a simplicial object in a category $\mathcal{C} $ as a contravariant functor $\boldsymbol{\Delta} ^{\mathrm{op}} \to \mathcal{C} $.
\end{remark}

\begin{construction}[Singular homology]\label{4.1.6}
	Let $X$ be a topological space. \Cref{4.1.1} gives the singular set $\Sing_{\bullet} (X)\in \sSet $. Recall that the free abelian group is a functor $\mathbb{Z} [-]: \Set  \to \Ab$ which maps sets to abelian groups generated by that set, and maps functions to the group homomorphism induced by that function on basis elements. Denote by $\mathrm{s} \Ab$ the category of simplicial abelian groups. We can extend the free abelian group functor to a functor $\mathbb{Z} [-]: \sSet  \to \mathrm{s} \Ab$ by applying it levelwise:
	\[
		\mathbb{Z} [\Sing(X)]_n=\mathbb{Z} [\Sing_n(X)],\qquad d_i^\mathbb{Z}  = \mathbb{Z} [d_i^{\Sing}]
	.\]
	We will denote by $C_n(X)$ the abelian group $\mathbb{Z} [\Sing_n(X)]$.

	Given the simplicial abelian group $C_{\bullet} (X)$, let $d_i : C_n \to C_{n-1} $ be the face map, and define the \emph{boundary morphism} on the basis simplices $\sigma\in\Sing_n(X)$ as
	\[\begin{tikzcd}[row sep=0em]
	{\partial_n:C_n} & {C_{n-1}} \\
	\sigma & {\displaystyle{\sum_{i=0}^n(-1)^id_i(\sigma)}}
	\arrow[from=1-1, to=1-2]
	\arrow[maps to, from=2-1, to=2-2]
	\end{tikzcd}\]
	and we extend it to all of $C_{\bullet} $ via linearity. We can offer a good description of what $d_i(\sigma)$ is defined as. On basis simplices, the face maps of the simplicial abelian group $C_{\bullet} (X)$ must agree with the face map $d_i : \Sing_n(X) \to \Sing_{n-1} (X)$ by the universal property of free abelian groups, and the face map $d_i$ in singular sets is precomposition by the coface map $d^i$ of geometric simplices. Thus, $d_i(\sigma)=\sigma\circ d^i\in\Sing_{n-1} (X)$ also holds between $C_n$ and $C_{n-1} $, so we will also write
	\[
		\partial_n : \sigma \mapsto \sum_{i=0}^{n} (-1)^i(\sigma\circ d^i)
	\]
	on basis simplices.

	We claim for now that $\partial _{n-1} \partial _n=0$. We prove this in \cref{4.1.?}, but the proof is lengthy, so we defer it for now. We thus have that $(C_{\bullet} (X),\partial _{\bullet} )$ is a chain complex of abelian groups. The assignment $C_{\bullet} (X)\mapsto (C_{\bullet} (X),\partial _{\bullet} )$ is functorial, which is explored in \cref{4.1.??}. We can thus compute the homology groups, typically denoted $H_{\bullet} (X)$ rather than $H_{\bullet} (C(X))$, which we recall are defined as
	\[
		H_n(X)=\frac{\Ker \partial _{n-1} }{\Im \partial _n} 
	.\]
	These homology groups are the \emph{singular homology} of the topological space $X$, and are the canonical and most useful construction of homology for topological spaces.
\end{construction}


\begin{remark}\label{4.1.7}
	\Cref{4.1.6} generalizes in a variety of ways. Firstly, we do not have to work with abelian groups. Given a (unital) ring $R$, we can define the $R$-valued homology $R$-modules $H_{\bullet} (X;R)$ by simply taking the free $R$-module on $\Sing_{\bullet} (X)$, rather than the free abelian group, and proceeding with \cref{4.1.6} in the same way. Of course, abelian groups are simply $\mathbb{Z} $-modules, so we have $H_{\bullet} (X;\mathbb{Z} )=H_{\bullet} (X)$. The other generalization is the notion of simplicial homology, which defines homology for a fairly broad class of simplicial sets known as \emph{simplicial complexes}. We will not introduce this notion in this paper in favor of exploring homotopical implicications of the functor $\Sing$, but the interested reader is encouraged to see \cref{at}, or to take a class on algebraic topology.
\end{remark}



\begin{proposition}\label{4.1.?}
	With notation as in \cref{4.1.6}, $\partial _{n-1} \partial _n=0$. In particular, $C_{\bullet} (X)$ is a chain complex.
\end{proposition}
\begin{proof}
	This proof is most straightforward when appealing to the simplicial identities of \cref{3.1.7}, rather than to the explicit description of simplicial operators in $C_{\bullet} (X)$. For this reason, let $d_i$ denote the $i$th face in the simplicial abelian group $C_{\bullet} (X)$, and let $\sigma\in C_n(X)$ be an $n$-simplex. We have
	\[\begin{tikzcd}
	\sigma & {\displaystyle{\sum_{i=0}^n(-1)^id_i(\sigma)}} & {\displaystyle{\sum_{j=0}^{n-1}(-1)^id_j\left(\sum_{i=0}^n(-1)^id_i(\sigma)\right)}.}
	\arrow["{\partial_n}", maps to, from=1-1, to=1-2]
	\arrow["{\partial_{n-1}}", maps to, from=1-2, to=1-3]
	\end{tikzcd}\]
	Since $d_j$ is a group homomorphism, it distributes over addition, giving
	\[
		\sum_{j=0}^{n-1} (-1)^jd_j\left( \sum_{i=0}^{n} (-1)^id_i(\sigma) \right) =\sum_{j=0}^{n-1} (-1)^j \sum_{i=0}^{n} d_j\left( (-1)^id_i(\sigma) \right) 
	.\]
	Note that $(-1)^ix$ is either $x$ or $-x$. For the former, $d_j((-1)^ix)=d_j(x)$, and for the latter, $d_j((-1)^ix)=d_j(-x)=-d_j(x)$, since $d_j$ is a group homomorphism. We may thus factor out the $(-1)^i$, giving
	\[
		\sum_{j=0}^{n-1}(-1)^j\sum_{i=0}^n d_j\left((-1)^id_i(\sigma)\right)=\sum_{j=0}^{n-1} (-1)^j \sum_{i=0}^{n} (-1)^id_jd_i(\sigma)=\sum_{j=0}^{n-1} \sum_{i=0}^{n} (-1)^j(-1)^i d_jd_i(\sigma)
	.\]
	a
\end{proof}


\begin{remark}\label{4.1.??}
	show functoriality
\end{remark}

